% Reduced extract from Kent E. Holsinger, Lecture Notes in Population Genetics.
% License: CC BY-SA 4.0. See SOURCES.md for attribution and source revision.
\documentclass{book}
\newcommand{\half}{\frac{1}{2}}
\newcommand{\fourth}{\frac{1}{4}}

\begin{document}
\chapter{The Hardy-Weinberg Principle and estimating allele frequencies}
\section*{The genetic composition of populations}

When I talk about the genetic composition of a population, I am referring to
three aspects of genetic variation within that population:\footnote{At each
locus I am talking about one locus at a time.}
\begin{enumerate}
\item The number of alleles at a locus.
\item The frequency of alleles at the locus.
\item The frequency of genotypes at the locus.
\end{enumerate}

It may not be immediately obvious why we need both genotype and allele
frequencies to describe the genetic composition of a population. Consider two
hypothetical populations:
\begin{center}
\begin{tabular}{lrrr}
             & $A_1A_1$ & $A_1A_2$ & $A_2A_2$ \\
\hline\hline
Population 1 &       50 &        0 &       50 \\
Population 2 &       25 &       50 &       25 \\
\hline
\end{tabular}
\end{center}

The frequency of $A_1$ is 0.5 in both populations, but the genotype
frequencies are very different. The following mating table uses the same
notation:
\begin{center}
\begin{tabular}{rcccc}
\hline\hline
                       &           & \multicolumn{3}{c}{Offspring genotype} \\
Female $\times$ Male   & Frequency & $A_1A_1$ & $A_1A_2$ & $A_2A_2$ \\
\hline
$A_1A_1 \times A_1A_1$ & $x_{11}^2$ & 1 & 0 & 0 \\
              $A_1A_2$ & $x_{11}x_{12}$ & $\half$ & $\half$ & 0 \\
              $A_2A_2$ & $x_{12}x_{22}$ & 0 & 1 & 0 \\
\hline
\end{tabular}
\end{center}

Taking these assumptions together allows us to calculate a frequency in the
pool of newly formed zygotes:
\[
\sum(\text{frequency of mating})(\text{frequency of genotype}) .
\]
\end{document}
